% ECONOMICS_PROBLEM: {"id": "H22-245", "title": "Corporate-tax incidence", "jel_code": "H22", "source_id": "OP-0474"}
% !TeX program = xelatex
\documentclass[11pt,a4paper]{article}
\usepackage[margin=23mm]{geometry}
\usepackage{fontspec}
\setmainfont{Latin Modern Roman}
\usepackage{amsmath,amssymb,mathtools}
\usepackage{xcolor,enumitem,booktabs,longtable,array,xurl,hyperref}
\definecolor{muted}{HTML}{53636C}
\hypersetup{colorlinks=true,urlcolor=blue,linkcolor=blue}
\setlength{\parindent}{0pt}
\setlength{\parskip}{5pt}
\newcommand{\R}{\mathbb R}
\newcommand{\E}{\mathbb E}
\newcommand{\N}{\mathbb N}
\newcommand{\ind}{\mathbf 1}
\newcommand{\Var}{\operatorname{Var}}
\newcommand{\Cov}{\operatorname{Cov}}
\newcommand{\supp}{\operatorname{supp}}
\DeclareMathOperator*{\argmin}{arg\,min}
\DeclareMathOperator*{\argmax}{arg\,max}
\newcommand{\entrymeta}[3]{{\sffamily\small\color{muted}\textbf{\texttt{#1}}\quad|\quad #2\par #3\par}}
\newcommand{\Problemref}[1]{\texttt{#1}}
\newcommand{\JELref}[2]{\texttt{#2} (JEL \texttt{#1})}
\begin{document}
\section*{Corporate-tax incidence}
\textbf{Problem ID:} \texttt{H22-245}\quad \textbf{JEL:} \texttt{H22}\par
% BEGIN ECONOMICS STATEMENT
\label{problem:OP-0474}
{\sffamily\small\color{muted}Original subject group: Taxation and social insurance\par}
\label{definitions:problem:30007756}
\textbf{Audit classification:} Published research agenda. The publisher verifies Auerbach’s 2006 unresolved-incidence discussion. It does not certify the exact extension remains unsolved. Fixed-price compensation and valid Shapley hybrid outcomes now replace ambiguous channel accounting.
\subsection*{Definitions and precise research target}
Fix a population with measure \(\mu\), mutually disjoint groups \(G_1,\ldots,G_K\), tax regimes \(p_0,p_1\), horizon \(H\), date-zero currency and a complete use of additional revenue. A regime includes its announcement date. Specify a structural equilibrium model or identification class for wages, prices, asset values and transfers, including ownership shares so corporate profits and rents are counted in household resources exactly once. Define \(U_i^H(p,x)\) as household expected utility when its budget receives the date-zero transfer \(x\), holding economy-wide equilibrium prices fixed for this individual compensation calculation. Assume continuity and strict monotonicity over an admissible transfer interval and that baseline utility lies in its range. The unique compensating variation \(CV_i\) solves
\[
 U_i^H(p_1,CV_i)=U_i^H(p_0,0).
\]
Positive \(CV_i\) is a loss under the reform. Let \(\Delta R_H\ne0\) be discounted incremental net public receipts and assume integrability. Define group incidence \(\iota_g=\int_{G_g}CV_i\,d\mu(i)/\Delta R_H\). Neither group shares nor loss-to-revenue ratios necessarily sum to one: distortion costs, insurance and chosen revenue use affect welfare. If compensation changes aggregate demand or prices, it is a separate compensated-equilibrium calculation and must be labelled accordingly.

For channel attribution, fix a finite set \(C\) of channels and, for each \(A\subseteq C\), a well-defined hybrid valuation \(v_g(A)\) with exactly those channels replaced by reform values. Such hybrids need not be equilibria; their construction is an additional attribution convention. The Shapley contribution of channel \(c\) is
\[
 \phi_{gc}=\sum_{A\subseteq C\setminus\{c\}}
 \frac{|A|!(|C|-|A|-1)!}{|C|!}\,[v_g(A\cup\{c\})-v_g(A)].
\]
Characterize or bound \(\iota_g\) and \(\phi_{gc}\) under stated capital-mobility, pricing, anticipation and data restrictions. The attribution identity is mathematical bookkeeping, not identification of causal channels. This is an editorial refinement of a published incidence question.
\subsection*{Variants and assumption changes}
\begin{enumerate}
\item \textbf{Announcement versus implementation (editorial).} Include the asset-price response on announcement in one horizon; start at implementation in another. Legacy shareholders' losses differ even under identical subsequent tax receipts.
\item \textbf{Closed versus open economy (editorial).} Hold capital supply fixed in a closed model or specify international asset supply and foreign ownership. Domestic group incidence and global total incidence are distinct targets.
\item \textbf{Transfer incidence versus full welfare (editorial).} Use quasilinear utility and fixed quantities as an accounting benchmark; compare with nonlinear utility and endogenous quantities using compensating variation.
\end{enumerate}
\subsection*{References and source audit}
Auerbach (2006), publisher abstract and journal metadata, Tax Policy and the Economy 20, 1--40: the incidence question and timing distinctions are verified. Full channel decomposition and robust empirical target are editorial extensions.
\begin{enumerate}\raggedright
\item\hypertarget{definitions:source:30007756:1}{} Alan J. Auerbach. 2006. \emph{Who Bears the Corporate Tax? A Review of What We Know}. Tax Policy and the Economy 20, 1–40; publisher abstract.
\url{https://www.journals.uchicago.edu/doi/10.1086/tpe.20.20061903}.
DOI: \url{https://doi.org/10.1086/tpe.20.20061903}.
\end{enumerate}

\subsection*{Common conventions: Source mathematical conventions}

These conventions apply to each entry unless explicitly replaced.

\textbf{Models and identification.} A primitive model \(m\) specifies all probability laws, feasible choices, information, equilibrium and measurement rules used by its target. The admissible class \(\mathcal M\) is fixed independently of outcomes. If \(P_O(m)\) is its observable probability law and \(\tau(m)\) the target, its identified set at observable law \(P\) is
\[
 \mathcal I_\tau(P)=\{\tau(m):m\in\mathcal M,\ P_O(m)=P\}.
\]
Point identification means that this set is a singleton. An empty set signals incompatibility of data and assumptions. Coordinate bounds need not describe the joint set. Bounds are sharp only if every retained target is attainable or arbitrarily approachable by compatible models. Finite bounds require explicit restrictions on unobserved quantities; unrestricted confounding, hidden wealth, extrapolation or arbitrary equilibrium selection can yield uninformative sets.

\textbf{Causal interventions.} A potential outcome \(Y_i(p)\) is the outcome under a complete policy regime \(p\), including timing, financing, information and market spillovers. The target population and units are fixed before treatment. With interference, use the complete assignment vector or a justified exposure mapping; individual no-interference notation alone cannot identify an economy-wide intervention. A difference of mean potential outcomes does not require the cross-world joint distribution. Conditioning on post-treatment migration, survival, enrollment or employment changes the estimand and needs separate justification.

\textbf{Statistical statements.} An estimator, confidence set, forecast or test specifies a sampling law, model class and loss/coverage criterion. Uniform \((1-\alpha)\)-coverage means \(\inf_{m\in\mathcal M}\Pr_m\{\tau(m)\in C_n\}\geq1-\alpha\), or an explicitly asymptotic version. The whole parameter space trivially has coverage, so informativeness requires a stated width, risk or power objective. A forecasting improvement is relative to a named benchmark, information set and evaluation population.

\textbf{Equilibrium and optimization.} An equilibrium comprises feasible plans, optimality, beliefs where needed, budget constraints and all market/resource clearing. The primitives or a stated benchmark define each correspondence \(\mathcal E(p;m)\). Nonemptiness is assumed or proved, never inferred just from compact policy sets. A selected-equilibrium target uses a declared selection rule; a robust target quantifies over all permitted equilibria. Use suprema or infima unless attainment is established. An \(\varepsilon\)-optimal policy is feasible and within \(\varepsilon\) of the finite optimum value. Report an empty feasible set separately.

\textbf{Welfare and accounting.} Normative utility, interpersonal normalization, social weights, numeraire, discounting and the use of public receipts are primitives. Behavioral utility need not coincide with the welfare criterion. Household welfare includes ownership income and rents once; adding profits again to compensating variation generally double-counts them. A monetary equivalent is defined by an explicit transfer and price convention, with monotonicity and range conditions. Average effects, mechanism contributions, transfers and welfare losses are different targets. Infinite sums require convergence and dynamic feasible plans require terminal or no-Ponzi conditions.

\textbf{Source versions.} A working-paper date, revision date and journal publication year are distinguished. A DOI for a journal version is not automatically the DOI of an earlier manuscript. Locators refer to the stated version and printed pages where specified. A metadata-only check does not verify a claim about a full-text conclusion. Every entry's source subsection narrows the provenance claim accordingly.
% END ECONOMICS STATEMENT
\end{document}
